\section{A complementary acceleration: coordinate blocks in normal-sector search}
\label{sec:sector-blocks}

The marked-interface results address what must be preserved after cutting.
This section concerns an earlier cost: finding normal-coordinate candidates
inside one supplied compatible quadrilateral sector.  It does not assume
that the sector is complete for disc discovery.  The contribution is an
exact decomposition of the matching search, a source-certified implementation,
and a genuine solid-torus family on which the maintained exhaustive
intersection iterator performs exponentially more candidate solves than the
decomposed iterator.  The linear-algebra principles are classical; the
source-level decomposition, its interaction with canonical triangle minima,
and the explicit geometric separation are the contributions developed here.

The comparison is with the repository snapshot identified earlier
\cite{ProveIt2026}.  Its current modules already contain support contraction,
matching-sign implications, grouped Euler minima, generic Q-corner screening,
and special algorithms at matching nullity at most three.  Reports 67 and 75
also retain anchor-LP and exact witness-reuse approaches.  None of those
results is being claimed again as a new theorem.  The present interface is
additive and opt-in; it does not change the recognizer's default dispatch.

\subsection{Matching-coordinate factors and canonical Euler characteristic}

Fix a source accepted by the native finite-manifold validator: a connected,
compact, orientable triangulated three-manifold with one torus boundary and
sphere or disc vertex links.  Let $T$ be its tetrahedron count.  A compatible
sector allows at most one quadrilateral type per tetrahedron.  After
eliminating triangle matching, write
\[
 L=\ker A\subseteq\mathbb R^k,\qquad
 C=L\cap\mathbb R_{\geq0}^k,
\]
where the $k$ coordinates are the allowed quadrilateral types.  A coordinate
block is a subset $I$ of these indices; $\mathbb R^I$ means vectors supported
on $I$.  All spaces and eliminations below are over the rationals when
implemented.

\begin{theorem}[The finest coordinate decomposition]
\label{sec:sector-coordinate-theorem}
There is a polynomial-time exact construction of pairwise disjoint blocks
$I_1,\ldots,I_s$ and a set $Z$ of coordinates identically zero on $L$ such that
\[
 L=\bigoplus_{j=1}^s L_j,\qquad
 L_j=L\cap\mathbb R^{I_j},\qquad
 C=\bigoplus_{j=1}^s C_j,\quad C_j=L_j\cap\mathbb R_{\geq0}^{I_j}.
\]
The nonzero blocks refine every other coordinate direct-sum decomposition.
Every extreme ray of $C$ belongs to exactly one nonzero $C_j$, and every
extreme ray of $C_j$ is an extreme ray of $C$.
\end{theorem}
\begin{proof}
Compute a canonical kernel basis $r_1,\ldots,r_d$ with free coordinates
$f_1,\ldots,f_d$ satisfying $(r_i)_{f_j}=\delta_{ij}$.  The native gauge uses
the rightmost independent free coordinates; any fixed canonical gauge works.
Join coordinates that occur together in the support of a basis row.  The
connected components of this graph on occurring coordinates are the $I_j$;
the nonoccurring coordinates form $Z$.  Each basis row is supported in one
component.  Grouping the basis rows therefore proves the direct-sum formula
for $L$, and coordinatewise nonnegativity proves the formula for $C$.

For the refinement assertion, consider any coordinate decomposition of $L$.
Project $r_i$ onto one of its summands.  Except in the summand containing
$f_i$, this projection has every free coordinate zero.  It must therefore
vanish, because restriction to the free coordinates is injective on $L$.
Thus every support of $r_i$ lies in a single summand, and so does every
connected component constructed above.  This argument also explains why
one must first canonicalize an arbitrarily mixed input basis.

A nonzero vector with nonzero components in two of the $C_j$ decomposes as
a sum of nonproportional feasible vectors and cannot generate an extreme
ray.  Conversely, if a vector supported in $I_j$ is decomposed into two
nonnegative vectors in $C$, both summands vanish outside $I_j$.  Its local
extremality therefore implies global extremality.  Rational elimination and
the support graph have polynomial bit complexity in their input encoding.
\end{proof}

This is an algebraic cone decomposition.  It does not claim that the
physical components of a normal surface are disjointly partitioned by these
blocks.  The earlier source-equation support decomposition in report 69 has
a different, geometric purpose.  Cancellation in eliminated matching rows
cannot by itself establish physical disjointness.

Let $X(q)$ be the canonical standard-coordinate lift: subtract the least
triangle potential at each global vertex.  Write $p_c(q)$ for its linear
corner potentials, $G_v$ for the corners at vertex $v$, and
$\lambda_v=\chi(\operatorname{link}v)$.  The exact source Euler compiler gives
\begin{equation}
\label{sec:sector-euler-formula}
 F(q):=\chi(X(q))=\gamma\cdot q
                  -\sum_v\lambda_v\min_{c\in G_v}p_c(q).
\end{equation}
Here $\lambda_v$ is one for a boundary vertex and two for an interior vertex.

\begin{lemma}[The canonical-lift defect]
\label{sec:sector-defect}
For $q,r\in C$, put
\[
 \delta_v(q,r)=\min_c\bigl(p_c(q)+p_c(r)\bigr)
                     -\min_c p_c(q)-\min_c p_c(r).
\]
Then $\delta_v(q,r)\geq0$ and
\[
 X(q)+X(r)=X(q+r)+\sum_v\delta_v(q,r)L_v,
\]
where $L_v$ is the standard-coordinate vertex-link vector.  Consequently,
if every $\lambda_v\geq0$, then
\[
 F(q+r)\leq F(q)+F(r),\qquad F(aq)=aF(q)\quad(a\geq0).
\]
\end{lemma}
\begin{proof}
Every sum $p_c(q)+p_c(r)$ is at least the sum of the two separate minima.
Subtracting the three canonical triangle formulas gives the displayed
coordinate identity; quadrilateral coordinates agree identically.
Applying the source Euler covector gives
$F(q+r)=F(q)+F(r)-\sum_v\lambda_v\delta_v(q,r)$.
The asserted inequality and positive homogeneity follow.
\end{proof}

\begin{theorem}[Exact block localization of the Euler maximum]
\label{sec:sector-euler-maximum}
Assume $C\ne\{0\}$ and every $\lambda_v\geq0$.  Omit blocks whose feasible
cone is zero.  Then
\[
 \max_{q\in C,\,\|q\|_1=1}F(q)
 =\max_j\ \max_{q\in C_j,\,\|q\|_1=1}F(q).
\]
In particular, the sector contains a canonical positive-Euler normal surface
if and only if one coordinate block does.  This remains true when triangle
minima couple several matching blocks.
\end{theorem}
\begin{proof}
The normalized sections are compact.  For a normalized $q$, decompose it
as $q=\sum_jq_j$ using Theorem~\ref{sec:sector-coordinate-theorem}.  Put
$w_j=\|q_j\|_1$ and, when $w_j>0$, $u_j=q_j/w_j$.
The preceding lemma gives
\[
 F(q)\leq\sum_jF(q_j)=\sum_{w_j>0}w_jF(u_j)
                         \leq\max_j\max F|_{C_j\cap\{\|q\|_1=1\}},
\]
since $\sum_jw_j=1$.  Every local normalized section is contained in the
global section, proving the reverse inequality.  Rational polyhedral
vertices suffice for the maxima of the convex piecewise-linear function;
clearing denominators produces integral normal coordinates.
\end{proof}

The sign condition is substantive.  For the abstract model
$C=\mathbb R_{\geq0}^2$ and $F(x,y)=\min(x,y)$, corresponding to a negative
minimum weight, both coordinate axes have value zero but the interior has
positive value.  This model is excluded by the source hypotheses.

Positive Euler is still not a disc predicate.  A positive ray may contain
a sphere or a boundary-inessential disc, and it need not be the ray required
by a downstream topological test.  Thus the theorem permits a complete
Euler exclusion or an Euler witness; it does not permit declaring an
essential-disc search exhausted after testing only the block Q rays.

\subsection{When complete standard rays also split}

For each global vertex $v$, choose an anchor $a_v\in G_v$.  Say that $v$
varies on a matching block $I_j$ if some restricted difference
$(p_c-p_{a_v})|_{L_j}$ is nonzero.  Connect any two blocks on which the same
vertex varies, and take connected components of this graph.  Call their
unions $J_1,\ldots,J_u$.  These definitions depend only on differences of
potentials, so they are invariant under additive linear changes of the
potential gauge.

\begin{theorem}[Standard rays and minimum-interaction clusters]
\label{sec:sector-standard-clusters}
The non-link standard extreme rays of the full sector are exactly the union
of the non-link standard extreme rays of its sectors restricted to
$J_1,\ldots,J_u$.  Every such ray occurs in one restricted sector only.
\end{theorem}
\begin{proof}
Write $q=\sum_iq_i$ according to these unions.  At a vertex $v$, all
nonconstant differences $p_c-p_{a_v}$ vary on at most one $J_i$.
After the common anchor term cancels in canonical normalization, the
triangle coordinates at $v$ therefore depend on $q_i$ alone; on every other
summand they are zero.  If no block varies at $v$, all canonical triangle
coordinates there are zero.  It follows coordinatewise that
\[
                         X(q)=\sum_i X(q_i).
\]

Every non-link standard extreme ray is canonical: a positive vertex-link
offset would split it into a nonzero canonical vector and a link vector.
The displayed identity then forces its quadrilaterals into one $J_i$.
Conversely, forbidding every Q coordinate outside $J_i$ defines a coordinate
face of the full standard cone.  An extreme ray of that face is extreme
in the full cone, and its canonical lift is the same source lift.
Different unions have disjoint Q supports, preventing duplicates.
\end{proof}

Let $k_j=|I_j|$, $d_j=\dim L_j$, and $h_j\leq k_j$ be the number of distinct
projected Q coordinate hyperplanes.  Complete Q enumeration now attempts
\begin{equation}
\label{sec:sector-block-cost}
                         \sum_j\binom{h_j}{d_j-1}
\end{equation}
candidate intersections, with zero-dimensional blocks omitted.  The same
argument handles feasible cones of smaller dimension than their matching
spaces.  Each candidate has polynomial bit cost; exact rational coefficients
are controlled by determinants of input-sized matrices.  Source preparation,
projection, actual coordinate output, and certificate replay must all be
charged.  For complete standard enumeration, the corresponding sum uses
the coarser $J_i$, with at most $O(|J_i|^2)$ potential hyperplanes per
cluster; the maintained low-dimensional envelope algorithms can be retained.

Consequently logarithmic \emph{maximum local} nullity gives a
quasi-polynomial supplied-sector bound even if total nullity is linear.
No theorem here bounds that maximum on arbitrary knot exteriors or supplies
a complete small family of sectors.  A single large cluster can retain the
original exponential cost.

\subsection{A genuine solid-torus family with exponential iterator separation}

The artifact constructs a simplicial solid torus by triangulating
$\Delta^2\times S^1$, with the circle divided into three edges.  Each
triangular prism has its three staircase tetrahedra, giving nine tetrahedra
and embedded boundary triangles.  Attach a sequence of triangulated
bipyramid balls along embedded boundary triangles.  Each attachment
preserves the solid torus.  The outgoing face in the construction contains
two newly introduced vertices and one existing vertex, so it is again an
embedded triangle suitable for the next attachment.

For the measured family, one ball has tetrahedra
\[
 (D,E,A,B),\qquad(D,E,B,C),\qquad(D,E,C,A),
\]
with cyclic internal face pairs.  Allow repository quadrilateral types
$(1,1,2)$ in those tetrahedra, and no quadrilaterals in the initial torus.

\begin{proposition}[Independent bipyramid matching factors]
\label{sec:sector-bipyramid}
With $b$ attached three-tetrahedron balls, the source has $T=9+3b$,
$k=3b$, and $d=2b$.  Its Q cone is the direct product of $b$ copies of
\[
 \{(u_0,u_1,u_2)\geq0:u_0+u_1=u_2\}.
\]
Each factor has the canonical matching basis
$(-1,1,0),(1,0,1)$ and its two extreme rays are
$(1,0,1)$ and $(0,1,1)$.  The global Q cone therefore has exactly $2b$ rays.
\end{proposition}
\begin{proof}
At the vertex $D$, the three internal face equations for triangle
potentials $z_0,z_1,z_2$ read
\[
 z_0+u_0=z_1,\qquad z_1+u_1=z_2+u_2,\qquad z_2=z_0.
\]
Their sum is $u_0+u_1=u_2$.  The cycle at $E$ gives its negative;
the equatorial vertex graphs are trees and supply no other Q equation.
Gluing the ball along a boundary triangle joins three pairs of distinct
vertex graphs by one edge each.  These are bridges, so they create no
additional triangle-potential cycle or Q relation.  This proves independence
inductively.  The claimed basis, dimension, and rays follow directly from
the displayed nonnegative equation.
\end{proof}

\begin{corollary}[Exponential-to-linear candidate separation]
\label{sec:sector-separation}
On this family the maintained exhaustive global Q iterator attempts exactly
\[
 \binom{3b}{2b-1}
       =\Theta\!\left(\frac{(27/4)^b}{\sqrt b}\right)
\]
candidate intersections.  The block iterator attempts exactly $3b$ and
returns the identical $2b$ rays.  Its complete construction and independent
block-coverage replay have polynomial bit complexity in $b$.
\end{corollary}
\begin{proof}
The three coordinate forms in each two-dimensional factor define three
distinct hyperplanes; forms on different factors are not proportional.
The global iterator explicitly visits every $(2b-1)$-subset of these $3b$
hyperplanes.  The new iterator visits the three one-hyperplane candidates
in each factor.  The kernel and output coordinates in this family have
polynomial encoding length, and the independent verifier makes only local
fixed-dimensional ray reconstructions after its polynomial global source
and splitting checks.  Stirling's formula gives the asymptotic count.
\end{proof}

This is a separation for the maintained iterator, not an exponential lower
bound for the geometric decision problem or for every ray-enumeration
algorithm.  In the triangular family, one local nonnegativity inequality
is redundant.  The same construction admits a stronger variant with a
four-tetrahedron bipyramid ball and types $(1,1,2,2)$: its relation is
$u_0+u_1=u_2+u_3$, its four coordinate inequalities are irredundant facets,
and it has four Q rays.  For $b$ such balls, the same proof gives
\[
 k=4b,\quad d=3b,\quad
 \text{global candidates}=\binom{4b}{3b-1},\quad
 \text{block candidates}=6b,\quad \text{Q rays}=4b.
\]
Thus removing redundant inequalities does not remove the general
independent-factor obstruction to global subset enumeration.

\subsection{A source-realized counterexample to splitting standard rays too far}

The two-ball triangular source already distinguishes the two theorems.
Put
\[
             q=(x,y,x+y,z,w,z+w),\qquad x,y,z,w\geq0.
\]
Its four Q rays are the four coordinate axes in these parameters.
The source's nonconstant vertex potential groups, after anchor shifts and
duplicate removal, have forms
\[
 \{0,y\}\ (\text{twice}),\quad \{0,-z\}\ (\text{twice}),
 \quad\{0,-x\},\quad\{0,w,w-x\}.
\]
Only the last group changes its minimizing form inside the Q cone; its
minimum is $\min(0,w-x)$.  The canonical standard fan is therefore the
nonnegative orthant subdivided by $x=w$.  It has five rays, adding
$e_x+e_w$ to the four coordinate axes.  The extra standard ray has
\[
                         q=(1,0,1,0,1,1).
\]
It is extreme because the independent conditions $y=0$, $z=0$, and $x=w$
leave one feasible dimension.  These exact source potential groups and the
five-ray list are also reconstructed by the independent dense verifier.
The minimum-interaction construction merges the two matching blocks, as
required.  A proposed implementation that enumerated standard rays in
the separate Q blocks would omit this genuine ray.

\subsection{Implementation, verification, and finite evidence}

The new producer is \texttt{fastunknot/sector\_blocks.py}.  It exposes
complete Q and standard enumeration, complete ray-set certificates, and a
canonical positive-Euler screen.  Its companion
\texttt{sector\_blocks\_verify.py} imports neither the block producer nor
its support-graph construction.  The checker reconstructs the dense
standard matching model from the original source.  For each submitted
block and each independently computed kernel basis vector, it checks that
the coordinate projection still satisfies all Q equations.  It checks
that every omitted coordinate vanishes on the full kernel.  These identities
prove the advertised direct-sum coverage without trusting a block count.
For standard certificates it additionally checks that each vertex varies
on at most one submitted cluster.  It then reconstructs every local ray
using the existing independent dense enumerator.

The positive-Euler certificate verifies matching, permitted types,
canonicality, and positive Euler directly; it asserts no component topology.
A negative Euler certificate includes complete local Q exhaustions and
the checked splitting.  Missing blocks, incorrect fine partitions,
missing or repeated rays, altered source identities, and noncanonical
positive vectors are rejected.  Candidate limits yield an inconclusive
result without an exhaustion certificate; callback interruption propagates.
Existing signed-hexadecimal integer transport is retained for ray entries.

The finite audit covers all 188 retained source sectors of matching
nullity four or five under the same support-generation policy used by the
maintained generic Euler tests.  All new Q ray sets equal the old complete
sets, and all Euler decisions equal the archived complete standard-coordinate
oracle.  Independent replay accepts every submitted certificate.
There are 180 split sectors, but only 19 have fewer candidate intersections;
the total drops from 1,275 to 1,110, with maximum individual factor $10/3$.
Euler outcomes remain 12 exclusions and 176 positive results.  These are
sector counts, not counts of distinct knot types or new recognition cases.

Nine focused test methods additionally cover dense changes of kernel basis,
zero cones, both enumeration phases on all 64 assignments of a three-tetrahedron
source, the genuine mixed standard ray, malicious certificate mutations,
integer transport, and cancellation.  The source family has also been
checked through eight triangular bipyramid attachments.  These finite
checks supplement the proofs rather than replacing their infinite claims.

\begin{center}
\begin{tabular}{r r r r r r r}
\hline
$b$ & $T$ & Old candidates & Block candidates & Old seconds & Block seconds & Paired gain\\
\hline
1 & 12 & 3 & 3 & 0.006644 & 0.006247 & 1.054\\
2 & 15 & 20 & 6 & 0.016219 & 0.013247 & 1.217\\
3 & 18 & 126 & 9 & 0.051434 & 0.021766 & 2.363\\
4 & 21 & 792 & 12 & 0.267364 & 0.032702 & 8.134\\
5 & 24 & 5005 & 15 & 1.848340 & 0.044701 & 39.974\\
\hline
\end{tabular}
\end{center}

The timing task includes complete source preparation, Q-ray production,
independent complete ray-set replay, and canonical ray-list hashing.
The new arm additionally serializes its block certificate.  Each row uses
three rounds with alternating arm order; times are arm medians and gains
are medians of paired ratios, which need not equal ratios of the displayed
medians.  Small cases are near parity and should not be read as strong
speed claims.  Larger cases exhibit the predicted avoided intersection
work.  The raw measurements and exact ray hashes are retained in
\texttt{sector\_block\_research/results/benchmark.json}; the 188-case audit
is beside it.  These are complete supplied-sector Q-enumeration measurements,
not end-to-end knot-recognition timings.

The next mathematical targets are more specific than a bound on total
nullity.  One can ask whether certified geometric reductions control the
largest surviving matching block; whether source-side forced-zero proofs
expose smaller blocks than the full linear kernel does; and whether a
small separator in the minimum-interaction graph admits an exact bounded
interface representation without expanding binary normal coordinates.
A separate boundary-aware theorem is required to preserve essential-disc
information when the Euler screen finds positive but topologically
unhelpful rays.  Those obligations remain open here.
